\documentclass[11pt]{article}
\usepackage[margin=1in]{geometry}
\usepackage{amsmath,amssymb}
\usepackage{enumitem}
\setlength{\parindent}{0pt}
\setlength{\parskip}{6pt}
\begin{document}
\section*{STAT 412 HW07 --- Question 8}

\subsection*{Problem}

\begin{quote}
A manufacturer claims that the average tensile strength of thread A exceeds the average tensile strength of thread B by at least 12 kilograms. Researchers wish to test this claim using a 0.1 level of significance. How large should the samples be if the power of the test is to be 0.99 when the true difference between thread types A and B is 6 kilograms? The population standard deviation for thread A is 6.28 kilograms and the population standard deviation for thread B is 5.68 kilograms.

Find the minimum sample size required. Round up to the nearest whole number.
\end{quote}

\subsection*{Solution}

Test H0: mu\_A-mu\_B >= 12 against a lower alternative. The distance to detect is 12-6=6 kg. With alpha=0.10 and power=0.99, n=[(z\_0.90+z\_0.99)\textasciicircum{}2(sigma\_A\textasciicircum{}2+sigma\_B\textasciicircum{}2)]/6\textasciicircum{}2 = [(1.2816+2.3263)\textasciicircum{}2(6.28\textasciicircum{}2+5.68\textasciicircum{}2)]/36 = 25.93, so round up to 26.

\subsection*{Final Answer}

\fbox{\begin{minipage}{0.94\linewidth}
Minimum sample size per group n=26.
\end{minipage}}

\end{document}
